\chapter{Implementation}

\section{Introduction}

\paragraph{}In order to be able to control the Baxter robot, a ROS environment has to be installed and configured on Ubuntu\footnote{\url{http://sdk.rethinkrobotics.com/wiki/Workstation_Setup}}. The Ubuntu Operating System is run on a virtual machine using the VMWare software. To test the inverse kinematics calculations described in Chapter \ref{chap:inv kin} without the physical Baxter robot, the Gazebo simulator software is used locally on the virtual machine. When the python code for the inverse kinematics is validated on the simulator, the same code can be tested on the physical robot. 

\section{Problems encountered and solutions}
Some problems were encountered during the project.\\
First the x-axis was oriented at 45$^\circ$ compared to the robot so an offset had to be added to have the x-axis if front of it, but it was to be positive for the right arm and negative for the left one. \\
The orientation of the axis had also to be verified. The group has indeed seen that the x-axis was in front of the robot after the off-set angle, both y-axis are pointing to the right and both z-axis are pointing to the bottom.\\
The sign of the value for the position in the inverse kinematic has to take it into account.\\
Another huge problem using this robot is the limit of the task and joint space of Baxter.
The task must be on a table so the range of displacement is quite limited but the place that was available on the table was close to the limit of its range. So the first point of reference had to be adjust to be at a good position.

\section{Python code}
\paragraph{}The code is written in Python programming language using rospy module\footnote{\url{https://github.com/ros/ros_comm}} and the baxter interface\footnote{\url{https://github.com/RethinkRobotics/baxter_interface}} to control the different parts of the robot. The code is made of 2 classes. 

\subsection{BaxterTask class}
The first one is the BaxterTask class that represents the task that we have to accomplish. To define the parameters of the task, a JSON file must be provided to the class containing everything needed such as the points the robot must achieve or the speed at which the robot must move. At initialisation of the task, all the points found in the JSON file are converted to joint angles using the inverse kinematics method explain in chapter~\ref{chap:inv kin}. The method \lstinline[language=Python, style=pythonColor]{inverse_kinematics} is in charge of this conversion. After the calculation of each point, the method \lstinline[language=Python, style=pythonColor]{check_angle_limits} verify that the angular positions asked by the inverse kinematics are physically feasible. If a value doesn't meet the requirements, a messaged is send as log with the missmatch in value.

\subsection{BaxterRobot class}
\paragraph{}The other class is the BaxterRobot class which manage all the communications with the robot. Finally, the main function boils down to the creation of an instance of each of these classes and calls the execute function of the BaxterTask class.


\section{Result}
 \begin{figure}[!ht]
	\centering
    \includegraphics[width = 0.7\textwidth]{Images/masterpiece-cropped.jpg}
    \caption{Final drawing of the robot}
    \label{fig:masterpiece}
\end{figure}